% Part: computability
% Chapter: recursive-functions
% Section: notations-pr-functions

\documentclass[../../../include/open-logic-section]{subfiles}

\begin{document}

\olfileid{cmp}{rec}{not}
\olsection{د بنسټيز بازګښت نښې}

د بنسټيزو بازګشتي تابعو د کره استقرايي بيان يوه ګټه دا ده چې دا
تابعې په منظم ډول بيانولے شو۔ د بېلګې په توګه، هرې داسې تابع ته
لاندې يوه «نښه» ټاکلے شو۔ د صفر، راتلونکي او پروجکشنونو دپاره په
ترتيب سره نښې~$\Zero$، $\Succ$ او~$\Proj{n}{i}$ وکاروئ۔ اوس فرض
کړئ~$h$ له يوې~$k$-ځاييزې تابع~$f$ او~$n$-ځاييزو تابعو~$g_0$،
\dots،~$g_{k-1}$ څخه په ترکيب تعريف شوې ده، او وروستيو تابعو ته
مو نښې~$F$، $G_0$، \dots،~$G_{k-1}$ ټاکلې دي۔ بيا د يوې نوې
نښې~$\fn{Comp}_{k,n}$ په کارولو سره، تابع~$h$ په
$\fn{Comp}_{k,n}[F,G_0,\dots,G_{k-1}]$ ښودلے شو۔

په بنسټيز بازګښت د تعريف شوو تابعو دپاره ورته نښې کارولے شو۔ فرض
کړئ~$(k+1)$-ځاييزه تابع~$h$ له~$k$-ځاييزې تابع~$f$ او
$(k+2)$-ځاييزې تابع~$g$ څخه په بنسټيز بازګښت تعريف شوې ده، او~$f$
او~$g$ ته ټاکل شوې نښې په ترتيب سره~$F$ او~$G$ دي۔ بيا~$h$ ته
ټاکل شوې نښه~$\fn{Rec}_k[F,G]$ ده۔

په ياد ولرئ چې د جمع تابع په بنسټيز بازګښت داسې تعريف شوې ده:
\begin{align*}
  \Add(x_0, 0) & = \Proj{1}{0}(x_0) = x_0\\
  \Add(x_0, y+1) & = \Succ(\Proj{3}{2}(x_0, y, \Add(x_0, y))) = \Add(x_0, y) +1
\end{align*}
دلته د~$f$ ونډه~$\Proj{1}{0}$ ترسره کوي، او د~$g$ ونډه
$\Succ(\Proj{3}{2}(x_0, y, z))$ ترسره کوي۔ وروستۍ تابع ته
$\fn{Comp}_{1,3}[\Succ,\Proj{3}{2}]$ نښه ټاکل شوې، ځکه دا له
يوې~$1$-ځاييزې تابع~$\Succ$ او~$3$-ځاييزې تابع~$\Proj{3}{2}$ څخه په
ترکيب د يوې تابع د تعريف پايله ده۔ په دې ترتيب، د جمع تابع په لاندې
نښه ښودلے شو:
\[
\fn{Rec}_1[\Proj{1}{0},\fn{Comp}_{1,3}[\Succ,\Proj{3}{2}]].
\]
د دې نښو لرل کله ناکله ګټور وي؛ د بېلګې په توګه، د بنسټيزو
بازګشتي تابعو د شمېرنې پر مهال۔

\begin{prob}
  د~$\Mult$ بشپړه بنسټيزه بازګشتي نښه ورکړئ۔
\end{prob}

\end{document}
